\documentclass[11pt]{article}
\usepackage{mathblatt}
\begin{document}
\blattkopf{Hypergeometrische Verteilung}{Lernblatt}
\verzeichniszeile{\verz{e1}{1 Genau k Treffer ohne Zurücklegen: die Situation erkennen (feste kleine Gesamtheit, Ziehen ohne Zurücklegen (Nr.~7–8)} \verztrenn \verz{e2}{2 Kumulieren gegen eine Schranke: hypergeometrische Einzelwahrscheinlichkeiten aufsummieren und die größte (oder kleinste) Trefferzahl gegen eine Schranke bestimmen (Nr.~9–10)} \verztrenn \verz{abhaken}{Das kann ich}}
% Einheit 1 von 2 – Genau k Treffer ohne Zurücklegen: die Situation erkennen (feste kleine Gesamtheit, Ziehen ohne Zurücklegen (bank/hypergeometrische-verteilung/e1.jsonl, zusammenbau v0.1)
%% TODO Zeitmarke Einheit 1: Marken-Zeile nicht lesbar
%% TODO Zweigzeile Teil 1: was hier gelernt wird (Satz in Schülersprache aus dem Einheitstitel) – Einheit 1
\einheitenkopf[e1]{Einheit 1 von 2 · Genau k Treffer ohne Zurücklegen: die Situation erkennen (feste kleine Gesamtheit, Ziehen ohne Zurücklegen}
\zweigzeile{Abitur GK}
\setcounter{aufgabe}{6}

%% TODO Titel: Ich-kann-Satz fehlt in der Bank (Platzhalter: Kettenname) – Nr. 7
%% TODO Anweisung: ein Satz über der ersten Teilaufgabe fehlt in der Bank – Nr. 7
\begin{aufgabe}{Genau k ohne Zurücklegen}
%% TODO \rechenplatz{n}: Schrittzahl des Grundfalls fehlt in der Bank (nur bei fester Schreibform, 2.3 b)
\begin{teile}
\teil Kreuze an, ohne zu rechnen: Aus einer Klasse mit $24$ Schülerinnen und Schülern werden $4$ für eine Umfrage ausgelost. Gefragt ist die Wahrscheinlichkeit für genau $2$ Mädchen. \\ \kreuz{ohne Zurücklegen – Auswahlen zählen} \\ \kreuz{mit Zurücklegen – Binomialmodell}
\teil In einer Urne liegen $5$ weiße und $3$ schwarze Kugeln. Es werden $3$ Kugeln ohne Zurücklegen gezogen. Berechne die Wahrscheinlichkeit, dass mindestens eine gezogene Kugel schwarz ist. \hfill (Abitur 2022 GK)
\teil Ein Verein hat $60$ Mitglieder, $9$ davon fahren E-Bike. $8$ Mitglieder werden zufällig für eine Befragung ausgewählt. Weise nach, dass die Wahrscheinlichkeit, dass genau $2$ davon E-Bike fahren, etwa $25{,}3\,\%$ beträgt. \hfill (Abitur 2024 LK)
\teil Von $20$ Bewerbungen kommen $8$ aus dem Ausland. $5$ Bewerbungen werden zufällig ausgewählt. Berechne die Wahrscheinlichkeit, dass höchstens $3$ davon aus dem Ausland kommen.
\teil An einem Kochkurs nehmen $120$ Personen teil, darunter Frau Keller. Unter den Teilnehmenden werden $4$ Gutscheine ausgelost, jede Person kann höchstens einen gewinnen. Berechne die Wahrscheinlichkeit, dass Frau Keller einen Gutschein gewinnt, und begründe, dass das Modell der Binomialverteilung dafür ungeeignet ist. \hfill (Abitur 2017 LK)
\end{teile}
\end{aufgabe}

%% TODO Titel: Ich-kann-Satz fehlt in der Bank (Platzhalter: Kettenname) – Nr. 8
%% TODO Anweisung: ein Satz über der ersten Teilaufgabe fehlt in der Bank – Nr. 8
\begin{aufgabe}{Genau k ohne Zurücklegen – Fehler finden, Begründen, Anwendung}
\begin{teile}
\teil Finde den Fehler und rechne richtig. Aus $20$ Losen mit $4$ Gewinnen werden $3$ gezogen. Mia rechnet für genau einen Gewinn: $3 \cdot 0{,}2 \cdot 0{,}8^2 = 0{,}384$.
\teil Begründe, warum sich beim Ziehen ohne Zurücklegen die Wahrscheinlichkeit für einen Treffer von Zug zu Zug ändert.
\teil Eine Lieferung enthält $50$ Bauteile, $3$ davon sind defekt. Der Kunde prüft $5$ zufällig entnommene Teile und nimmt die Lieferung nur an, wenn keines defekt ist. Mit welcher Wahrscheinlichkeit lehnt er die Lieferung ab?
\end{teile}
\end{aufgabe}

\clearpage
% Einheit 2 von 2 – Kumulieren gegen eine Schranke: hypergeometrische Einzelwahrscheinlichkeiten aufsummieren und die größte (oder kleinste) Trefferzahl gegen eine Schranke bestimmen (bank/hypergeometrische-verteilung/e2.jsonl, zusammenbau v0.1)
%% TODO Zeitmarke Einheit 2: Marken-Zeile nicht lesbar
%% TODO Zweigzeile Teil 1: was hier gelernt wird (Satz in Schülersprache aus dem Einheitstitel) – Einheit 2
\einheitenkopf[e2]{Einheit 2 von 2 · Kumulieren gegen eine Schranke: hypergeometrische Einzelwahrscheinlichkeiten aufsummieren und die größte (oder kleinste) Trefferzahl gegen eine Schranke bestimmen}
\zweigzeile{Abitur LK}
\setcounter{aufgabe}{8}

%% TODO Titel: Ich-kann-Satz fehlt in der Bank (Platzhalter: Kettenname) – Nr. 9
%% TODO Anweisung: ein Satz über der ersten Teilaufgabe fehlt in der Bank – Nr. 9
\begin{aufgabe}{Kumulieren}
%% TODO \rechenplatz{n}: Schrittzahl des Grundfalls fehlt in der Bank (nur bei fester Schreibform, 2.3 b)
\begin{teile}
\teil Kreuze den richtigen Ansatz an, ohne zu rechnen: Aus $10$ Kugeln, $4$ davon rot, werden $3$ ohne Zurücklegen gezogen. Gesucht ist $P(X = 2)$ für $X$ = Zahl der roten Kugeln. \\ \kreuz{$\frac{(4 \text{ über } 2) \cdot (6 \text{ über } 1)}{(10 \text{ über } 3)}$} \\ \kreuz{$(3 \text{ über } 2) \cdot 0{,}4^2 \cdot 0{,}6$} \\ \kreuz{$\frac{(4 \text{ über } 2)}{(10 \text{ über } 3)}$}
\teil Aus einer Urne mit $10$ Kugeln, davon $4$ rot, werden $3$ ohne Zurücklegen gezogen; $X$ ist die Zahl der roten Kugeln. Bestimme die größte Zahl $k$, für die $P(X \le k) < 0{,}7$ gilt, und belege beide Nachbarwerte.
\teil Eine Reisegruppe besteht aus $24$ Personen, $9$ davon sprechen Spanisch. Für einen Ausflug werden $6$ Personen zufällig ausgewählt; $X$ ist die Zahl der Spanischsprechenden. Ermittle das größtmögliche $n$ mit $P(X \le n) < 40\,\%$. \hfill (Abitur 2023 LK)
\end{teile}
\end{aufgabe}

%% TODO Titel: Ich-kann-Satz fehlt in der Bank (Platzhalter: Kettenname) – Nr. 10
%% TODO Anweisung: ein Satz über der ersten Teilaufgabe fehlt in der Bank – Nr. 10
\begin{aufgabe}{Kumulieren – Fehler finden, Begründen}
\begin{teile}
\teil Finde den Fehler und rechne richtig. Von $10$ Glühbirnen sind $4$ defekt; $3$ werden entnommen, $X$ zählt die defekten. Gesucht ist das größte $k$ mit $P(X \le k) < 0{,}66$. Nina rechnet mit $p = 0{,}4$: $P(X \le 1) = 0{,}6^3 + 3 \cdot 0{,}4 \cdot 0{,}6^2 = 0{,}648 < 0{,}66$, also $k = 1$.
\teil Begründe, warum man für die größte Trefferzahl unter einer Schranke zwei kumulierte Werte angeben muss.
\end{teile}
\end{aufgabe}

% Abhakseite „Das kann ich“ – Zone nur im Gesamt (\mitzone)
%% TODO Ich-kann-Sätze: Platzhalter wie in den Aufgabendateien
\begin{abhakseite}
\ifdefined\mitzone
\abhakgruppe{Kennst du schon}
\abhak{1}{Binomialkoeffizient (n über k) als Anzahl der Auswahlen, Rechnertaste und kleine Werte im Kopf}
\abhak{2}{Pfadregeln ohne Zurücklegen (Bruchkette mit schrumpfendem Nenner) und das Gegenereignis}
\abhak{3}{Das Binomialmodell und seine Bedingungen (zur Abgrenzung: wann es ungeeignet ist)}
\abhak{4}{Brüche multiplizieren und kürzen}
\abhak{5}{Pfadregeln ohne Zurücklegen (Bruchkette mit schrumpfendem Nenner) und das Gegenereignis – Fehler finden}
\abhak{6}{Pfadregeln ohne Zurücklegen (Bruchkette mit schrumpfendem Nenner) und das Gegenereignis – selbst rechnen}
\fi
\abhakgruppe{Einheit 1 · Genau k Treffer ohne Zurücklegen: die Situation erkennen (feste kleine Gesamtheit, Ziehen ohne Zurücklegen}
\abhak{7}{Genau k ohne Zurücklegen}
\abhak{8}{Genau k ohne Zurücklegen – Fehler finden, Begründen, Anwendung}
\abhakgruppe{Einheit 2 · Kumulieren gegen eine Schranke: hypergeometrische Einzelwahrscheinlichkeiten aufsummieren und die größte (oder kleinste) Trefferzahl gegen eine Schranke bestimmen}
\abhak{9}{Kumulieren}
\abhak{10}{Kumulieren – Fehler finden, Begründen}
\end{abhakseite}

\end{document}
